\chapter{Editing Structures}
\label{ch:editing}

Every operation in this chapter is a single undo step.
\menu{Edit}\menusep\menu{Undo} (\kbd{Ctrl}+\kbd{Z}) and
\menu{Edit}\menusep\menu{Redo} (\kbd{Ctrl}+\kbd{Shift}+\kbd{Z}) walk a
per-document history up to 50 states deep.

\section{Atoms}

\begin{description}
  \item[\menu{Edit}\menusep\menu{Add Atom}] \kbd{Ctrl}+\kbd{Shift}+\kbd{A}.
    Choose the element from the graphical periodic table and type exact
    $x$, $y$, $z$ coordinates. For placing atoms by eye instead, use
    Insertion mode in the viewport (\kbd{I}).
  \item[\menu{Edit}\menusep\menu{Selection}\menusep\menu{Change Element of Selection}]
    Replaces the species of every selected atom, again via the periodic
    table.
  \item[\menu{Edit}\menusep\menu{Selection}\menusep\menu{Translate Selection}]
    Shifts the selection by a $\Delta x$, $\Delta y$, $\Delta z$ vector in
    \AA.
  \item[\menu{Edit}\menusep\menu{Delete Selected Atoms}] \kbd{Delete}.
    Removes the selection; with the viewport focused in Selection mode,
    \kbd{Backspace} does the same.
\end{description}

Deleting atoms re-indexes everything that referenced them --- bond
overrides, bond orders and per-atom fields all follow correctly.

\section{The periodic table selector}

Wherever an element must be chosen, Calango opens a graphical periodic
table: $Z = 1$ to 118 in the standard 18-group layout, f-block below,
coloured by chemical family. One click selects and closes.

\section{Bond perception}
\label{sec:bonds}

By default Calango perceives bonds by distance: atoms $i$ and $j$ are
bonded when
\[
  d_{ij} < t \cdot \bigl(r_\text{cov}(i) + r_\text{cov}(j)\bigr),
\]
with covalent radii from Cordero \emph{et al.} and a tolerance $t$ you
control. With a periodic cell the minimum-image convention applies, so
bonds across cell boundaries are found and drawn correctly.

\menu{Edit}\menusep\menu{Bond Editor} (\kbd{Ctrl}+\kbd{B}) exposes both the
automatic rule and manual overrides:

\begin{description}
  \item[\ui{Automatic bond detection}] Turn it off to draw \emph{only}
    manually added bonds --- the right choice for coarse-grained or
    schematic figures.
  \item[\ui{Covalent cutoff multiplier}] The tolerance $t$, 0.50--2.50,
    default 1.15. Applied live, so you can dial it until the picture is
    right.
  \item[\ui{Manual Bonds}] Two 1-based atom-index spin boxes, or
    \ui{From Selection} to fill them from a two-atom viewport selection.
    The info line shows the current pair's distance next to the automatic
    cutoff for that pair, which makes it obvious whether a bond is
    borderline.
  \item[\ui{Add Bond}] Forces a bond regardless of distance.
  \item[\ui{Suppress Bond}] Removes an automatically perceived bond.
  \item[\ui{Active overrides}] Lists every override as \emph{added} or
    \emph{suppressed}, with the pair and its distance.
    \ui{Clear Override} and \ui{Clear All} undo them.
\end{description}

Overrides live on the structure and are saved in the project file.

\section{Bond orders}
\label{sec:bond-orders}

Calango never guesses multiple bonds. Every perceived bond starts as a
single bond, and orders are assigned deliberately --- distance-based
perception of double and triple bonds is unreliable, particularly for
inorganic and metallic systems where short contacts are common.

To assign one, select exactly two atoms in the viewport
(\kbd{Ctrl}+click the second) and use the \ui{Bond order} buttons in the
\ui{Representation} panel: \ui{Single}, \ui{Double}, \ui{Triple}. The
buttons are disabled until a pair is selected.

An order of two or three renders as that many parallel cylinders, and also
forces the bond to exist even if the pair lies outside the automatic
cutoff. Orders persist in project files.

\begin{calnote}
A consequence worth knowing: molecules like benzene or CO$_2$ render with
all-single bonds when first opened. Assign the orders you want to display;
they are a presentation property, not an input to any calculation.
\end{calnote}
